\section{Diskussion}
\label{sec:Diskussion}

Ziel des Versuches war die Bestimmung der effektiven Masse von Elektronen in
einem Halbleiter durch die Auswertung der Faraday-Rotation. Die effektiven Massen
wurden zu
\begin{align*}
    m^*_{1.2dot} &= \qty{7.0+-1.2e-32}{kg} = \qty{0.07+-0.013}{m_e} \\
    m^*_{2.8dot} &= \qty{1.2+-0.5e-31}{kg} = \qty{0.13+-0.05}{m_e}
\end{align*}
bestimmt. Diese Werte liegen deutlich unter der Ruhemasse eines freien Elektrons,
was auf die Wechselwirkung der Elektronen mit dem Kristallgitter zurückzuführen ist, 
die zu einer Verringerung der effektiven Masse führt. Für den Literaturwert der
effektiven Masse von Elektronen in GaAs wird ein Wert von
\begin{equation*}
    m^*_\text{GaAs} = 0.067 \cdot m_e
\end{equation*}
angegeben \cite{effmasse}. Für Probe 1 liegt das im Fehlerbereich, bei Probe 2 
hingegen liegt der ermittelte Wert mindestens 16.25 \% darüber. Angesichts der 
Tatsache, dass sich einerseits eine so präzise Messung ergibt, liegt es nahe, dass 
im Experiment zu 2. Probe ein systematischer Fehler vorliegt. Es kann sein, dass das 
höher dotierte GaAs ggf. durchlässiger ist, was zu einer geringeren Absorption der
Strahlung führt, wodurch die Messung der Faraday-Rotation überschätzt wird. Andererseits 
wäre die manuelle Nullabstimmung am Oszilloskop eine mögliche Fehlerquelle. Diese 
läuft durch Abschätzen mit dem bloßen Auge sowie Einstellen per Hand am Goniometer 
ab. Bei potenziellem Rauschen erschwert sich die genaue Bestimmung der Nullposition,
was in einer fehlerhaften Detektion der gemessenen Rotationswinkel führt. Das zeigt 
sich zum Teil auch in \autoref{fig:faradayrotation}, wo jeweils der dritte und 
vierte Messwert für die Auswertung nicht verwertbar sind. Diese beiden Ausreißer 
könnten jedoch auch durch die Interferenzfilter bedingt sein, da es teilweise 
vorkam, dass sich das Signal am Oszilloskop kontinuierlich verstärkt. Die restlichen 
Filter lieferten jedoch verhältnismäßig stabile Signale, sodass der Großteil der
Messwerte nicht besonders vom Fit abweicht.
\\
Abschließend bleibt festzuhalten, dass der Versuch das Verständnis für den 
Zusammenhang zwischen Magnetfeld, Lichtpolarisation und Bandstruktur in Halbleitern
erfolgreich vertieft hat. Die bestimmten effektiven Massen bestätigen den
Literaturwert für GaAs im Rahmen der experimentellen Genauigkeit und zeigen 
gleichzeitig die Grenzen des einfachen Modells bei höheren Dotierungen auf.